% Figure 1 — three-spindle puzzle (Problem 2-26).
% Semantic reconstruction of figures/scans/calcu5-fig1.png.
% Oblique board with mild perspective (front a little wider than the back,
% left edge more diagonal, right edge steeper). Four disks + spindle;
% two empty pegs. Labels and caption are TeX.

\begin{tikzpicture}[reconstruction, line cap=round, line join=round,
  line width=0.9pt, x=1mm, y=1mm]

  % Board corners — match the scan's nearly-horizontal front/back.
  \coordinate (FL) at (0.0, 2.0);
  \coordinate (FR) at (106.0, 6.0);
  \coordinate (BL) at (18.0, 50.0);
  \coordinate (BR) at (116.0, 52.5);

  % Peg 2 sits on the back edge (scan: the edge is interrupted, not drawn through).
  \def\pxtwo{70.0}
  \def\pytwo{51.3}
  \def\prad{5.5}
  \def\pry{2.1}

  \draw (FL) -- (FR) -- (BR);
  \draw (FL) -- (BL);
  \draw (BL) -- ({\pxtwo-\prad}, {\pytwo});
  \draw ({\pxtwo+\prad}, {\pytwo}) -- (BR);

  % Empty peg: opaque white, two ellipses, two generators.
  \newcommand{\hanoipeg}[5]{% #1,#2 base centre; #3 rx; #4 ry; #5 height
    \fill[white]
      ({#1-#3}, {#2}) -- ({#1-#3}, {#2+#5})
      arc[start angle=180, end angle=0, x radius=#3, y radius=#4]
      -- ({#1+#3}, {#2})
      arc[start angle=0, end angle=-180, x radius=#3, y radius=#4]
      -- cycle;
    \fill[white] ({#1}, {#2}) ellipse [x radius=#3, y radius=#4];
    \fill[white] ({#1}, {#2+#5}) ellipse [x radius=#3, y radius=#4];
    \draw ({#1}, {#2}) ellipse [x radius=#3, y radius=#4];
    \draw ({#1-#3}, {#2}) -- ({#1-#3}, {#2+#5});
    \draw ({#1+#3}, {#2}) -- ({#1+#3}, {#2+#5});
    \draw ({#1}, {#2+#5}) ellipse [x radius=#3, y radius=#4];
  }

  % Disk: shaded cylindrical wall (light from the left, as in the scan),
  % white top, full bottom ellipse.
  \newcommand{\hanoidisk}[5]{% #1,#2 top centre; #3 rx; #4 ry; #5 thickness
    \shade[left color=black!18, right color=black!38]
      ({#1-#3}, {#2-#5}) -- ({#1-#3}, {#2})
      arc[start angle=180, end angle=360, x radius=#3, y radius=#4]
      -- ({#1+#3}, {#2-#5})
      arc[start angle=360, end angle=180, x radius=#3, y radius=#4]
      -- cycle;
    \draw ({#1-#3}, {#2}) -- ({#1-#3}, {#2-#5});
    \draw ({#1+#3}, {#2}) -- ({#1+#3}, {#2-#5});
    \draw ({#1}, {#2-#5}) ellipse [x radius=#3, y radius=#4];
    \fill[white] ({#1}, {#2}) ellipse [x radius=#3, y radius=#4];
    \draw ({#1}, {#2}) ellipse [x radius=#3, y radius=#4];
  }

  % Peg 2 — on the rear edge, tall. Peg 3 is further right and forward;
  % the scan leaves a clear gap between them.
  \hanoipeg{\pxtwo}{\pytwo}{\prad}{\pry}{31}
  \node[font=\large] at (\pxtwo, 41.5) {2};

  \hanoipeg{98.0}{18.5}{\prad}{\pry}{28}
  \node[font=\large] at (98.0, 9.5) {3};

  % Stack, bottom disk first. Four rings + a short spindle (same radius as the pegs).
  \hanoidisk{37}{24.0}{21.5}{8.0}{4.4}
  \hanoidisk{37}{28.0}{17.4}{6.5}{4.0}
  \hanoidisk{37}{31.6}{13.6}{5.05}{3.6}
  \hanoidisk{37}{34.8}{10.0}{3.7}{3.2}
  \hanoipeg{37}{34.8}{5.5}{2.1}{6.8}

  \node[anchor=west, font=\normalsize, inner sep=0pt]
    at (0, -3.0) {\textls[260]{FIGURE 1}};
\end{tikzpicture}
